\section{Программные интерфейсы и интеграция}

\subsection{C++ API}

Публичный C++ API состоит из трёх уровней:

\paragraph{Уровень устройства.}
\texttt{Device} и \texttt{Stream} управляют выбором вычислительного устройства
и очередью операций. Пользователь создаёт устройство и поток один раз при запуске.

\paragraph{Уровень сцены.}
Объекты-хранилища (\texttt{BodyStore}, \texttt{ShapeStore}, \texttt{JointStore})
принимают описательные структуры (\texttt{BodyParams}, \texttt{ShapeParams},
\texttt{JointParams}) и управляют памятью на устройстве.
После заполнения вызывается \texttt{upload()}, переносящий данные на графический процессор.

\paragraph{Уровень симуляции.}
\texttt{Broadphase}, \texttt{Narrowphase}, \texttt{XpbdSolver}
и \texttt{RayQuery} реализуют этапы шага симуляции.
Цикл симуляции описан в листинге~\ref{alg:xpbd} в разделе~4.1.

\subsection{Python-биндинг через nanobind}

Модуль \texttt{dyphur\_py} предоставляет Python-интерфейс к основным компонентам движка через библиотеку 
\textbf{nanobind}. Nanobind обеспечивает нулевое копирование данных при работе с массивами numpy: устройство 
предоставляет указатель через интерфейс DLPack/CUDA Array Interface, не перенося данные на центральный процессор.

Контактный сенсор (\texttt{ContactSensor}) реализован как отдельный компонент с подпиской на контактные 
события по идентификаторам тел. Тест \texttt{test\_contact\_sensor.py} верифицирует его работу: 1024 кубов 
падают на статическую площадку, после чего проверяется, что каждый куб зарегистрировал хотя бы один контакт.

\subsection{Инструмент визуализации}

Для визуального контроля результатов симуляции реализован автономный инструмент \texttt{tools/viz/} на базе 
библиотеки \textbf{Magnum}.

Инструмент поддерживает три режима работы:

\begin{itemize}
  \item \textit{snapshot}~--- безголовый рендеринг через EGL (без монитора),
    вывод PNG-изображения; применяется для иллюстраций в документации.
  \item \textit{replay}~--- интерактивное окно GLFW с орбит-камерой;
    управление: перетаскивание мышью, прокрутка, пробел (пауза), стрелки (шаг по кадрам).
  \item \textit{live}~--- чтение из разделяемой памяти POSIX в реальном времени
    (на стороне движка~--- компонент \texttt{VizSink}).
\end{itemize}

Каждый демо-исполняемый файл записывает файл \texttt{.scene} (каталог форм тел) вместе с \texttt{.trajectory}, 
что позволяет визуализатору восстановить геометрию из данных поз.

\subsection{Система сборки}

Сборка проекта управляется CMake~3.25+ с генератором \textbf{Ninja Multi-Config}. Зависимости (Eigen, Catch2, 
spdlog, fmt, pugixml, nanobind) управляются через \textbf{vcpkg} в manifest-режиме с базовым SHA-пином 
(\texttt{vcpkg-configuration.json}). AdaptiveCpp устанавливается отдельно как системный пакет.

Предопределённые workflow-пресеты:

\begin{itemize}
  \item \texttt{dev}~--- Debug-сборка с CUDA, конфигурация для clangd (IDE);
  \item \texttt{ci-linux-cuda}~--- полный CI на NVIDIA;
  \item \texttt{ci-linux-hip}~--- полный CI на AMD;
  \item \texttt{ci-linux-cpu}~--- бэкенд ЦП с матрицей санитайзеров ASAN/UBSan/TSan;
  \item \texttt{ci-determinism}~--- только Release + golden trace.
\end{itemize}

Выбор пресета полностью определяет бэкенд, режим сборки и набор тестов. Разработчику достаточно одной команды: 
\texttt{cmake -{}-workflow -{}-preset=dev}.

\subsection{Выводы по разделу}

C++ API реализует принцип единственной ответственности: каждый компонент отвечает за один аспект~--- владение 
памятью, вычисление широкой фазы, решение ограничений~--- и взаимодействует с остальными только через явные 
представления-View. Python-интерфейс открывает движок для применения в задачах машинного обучения без 
дополнительных накладных расходов на копирование данных.
